\chapter*{Tóm tắt}
\label{summary}
\par Trong bối cảnh chuyển đổi số trong lĩnh vực y tế ngày càng phát triển, nhu cầu ứng dụng công nghệ thông tin vào việc theo dõi và quản lý sức khỏe từ xa trở nên ngày càng cần thiết. Đặc biệt, các bệnh mãn tính như tăng huyết áp và đái tháo đường đang ngày càng phổ biến, đòi hỏi người bệnh phải thường xuyên ghi nhận và kiểm soát các chỉ số sức khỏe quan trọng như huyết áp, nhịp tim, đường huyết và một số chỉ số sinh tồn cơ bản khác. Tuy nhiên, trong thực tế, việc theo dõi bệnh nhân thường chỉ được thực hiện thông qua các lần tái khám định kỳ tại cơ sở y tế, khiến bác sĩ khó nắm bắt đầy đủ diễn biến sức khỏe của bệnh nhân trong đời sống hằng ngày. Bên cạnh đó, việc ghi chép thủ công các chỉ số sức khỏe cũng dễ xảy ra sai sót, thất lạc dữ liệu và gây khó khăn trong quá trình tổng hợp, đánh giá tình trạng bệnh.

\par Từ những vấn đề trên, nhóm đã nghiên cứu và xây dựng một hệ thống theo dõi bệnh nhân từ xa dành cho người bệnh tăng huyết áp và đái tháo đường. Hệ thống bao gồm ứng dụng di động dành cho bệnh nhân, ứng dụng web dành cho bác sĩ, trang quản trị dành cho quản trị viên và các chức năng hỗ trợ nhân viên y tế trong quá trình theo dõi bệnh nhân. Thông qua ứng dụng di động, bệnh nhân có thể nhập các chỉ số sức khỏe hằng ngày, xem lại lịch sử đo, theo dõi cảnh báo và nhận các nhắc nhở cần thiết. Về phía bác sĩ, hệ thống web cho phép quản lý danh sách bệnh nhân, xem chi tiết lịch sử đo, quan sát xu hướng thay đổi của các chỉ số thông qua biểu đồ trực quan, cấu hình ngưỡng an toàn và tiếp nhận cảnh báo khi dữ liệu vượt ngưỡng. Ngoài ra, hệ thống còn hỗ trợ thông báo, nhắc nhở và trao đổi giữa bệnh nhân với nhân viên y tế, góp phần nâng cao tính liên tục trong quá trình theo dõi sức khỏe.

\par Nhóm hy vọng hệ thống sau khi được hoàn thiện có thể góp phần minh họa tính khả thi của mô hình Remote Patient Monitoring trong việc hỗ trợ quản lý bệnh mãn tính. Hệ thống không nhằm thay thế hoạt động khám chữa bệnh trực tiếp, mà đóng vai trò như một công cụ hỗ trợ giúp bệnh nhân chủ động hơn trong việc theo dõi sức khỏe, đồng thời giúp bác sĩ và nhân viên y tế có thêm dữ liệu tham khảo để phát hiện sớm các dấu hiệu bất thường. Trong tương lai, hệ thống có thể tiếp tục được mở rộng theo hướng tích hợp trực tiếp với các thiết bị đo sức khỏe, bổ sung các thuật toán phân tích dữ liệu và triển khai thực nghiệm trong môi trường y tế thực tế nhằm mang lại hiệu quả cao hơn cho công tác chăm sóc sức khỏe từ xa.